\documentclass[11pt,a4paper]{article}

\usepackage[utf8]{inputenc}
\usepackage[T1]{fontenc}
\usepackage[a4paper,margin=2.5cm]{geometry}
\usepackage{amsmath,amssymb,amsthm}
\usepackage{mathtools}
\usepackage{booktabs}
\usepackage{enumitem}
\usepackage[expansion=false]{microtype}
\usepackage[hidelinks]{hyperref}
\usepackage[numbers,sort&compress]{natbib}
% This document uses natbib + bibtex for maximum portability.
% For a biblatex/biber setup, delete the natbib line above and the
% \bibliographystyle/\bibliography lines at the end, then use instead:
%   \usepackage[backend=biber,style=numeric-comp,sorting=none,maxbibnames=99]{biblatex}
%   \addbibresource{references.bib}
% and replace \citet{...}/\citep{...} with \textcite{...}/\autocite{...},
% and the two bibliography commands at the end with
% \printbibliography[title={References}].

% --- Theorem-style environments ----------------------------------------------
\theoremstyle{definition}
\newtheorem{condition}{Condition}
\newtheorem{target}{Target}

% --- Notation shortcuts ------------------------------------------------------
\newcommand{\Voc}{\mathcal{V}}
\newcommand{\Fclass}{\mathcal{F}}
\newcommand{\RR}{\mathbb{R}}
\newcommand{\EE}{\mathbb{E}}
\newcommand{\Hent}{\mathcal{H}}

\title{\textbf{Compressed Sensing of LLMs}\\[0.3em]
\large Query-Efficient Recovery of a Black-Box Next-Token Function}
\author{Master's Thesis Proposal\\[0.4em]
\normalsize MSc in Computer Science / Machine Learning, Data Science and Artificial Intelligence\\
\normalsize 30 ECTS \quad$\cdot$\quad Supervisor: Prof.\ A.\ Jung\\
\normalsize Department of Computer Science, Aalto University}
\date{}

\begin{document}
\maketitle

% =============================================================================
\section{The Abstraction}

We model a language model as a single deterministic nonlinear function
\[
  f_\theta : \Voc^L \to \RR^{\Voc}, \qquad c \mapsto z = f_\theta(c),
\]
mapping a context window $c$ (a token sequence of length up to $L$) to the logit
vector $z$ over the next token, where $\Voc$ is the vocabulary. 

Auditing a closed LLM can be framed as a function recovery problem: Given observed 
conversations, estimate the underlying function $f_\theta$.

The thesis asks the compressed-sensing question in its purest form: 
treating $f_\theta$ as an unknown high-dimensional signal, how much of 
it can be recovered from $m \ll |\Voc|^L$ evaluations, under what 
structural conditions, and with what query designs?

There are no separate ``attacks,'' only different \emph{targets} --- the whole
function, a restriction, or a functional of it --- recovered through one
evaluation operator.

% =============================================================================
\section{Why Compressed Sensing Applies: the Compressibility Priors}

A generic function $\Voc^L \to \RR^{\Voc}$ is unrecoverable from few samples ---
the domain has size $|\Voc|^L$. Recovery is possible only because $f_\theta$ is
\emph{not} generic. Four structural facts play the role that sparsity plays in
classical compressed sensing~\citep{candes2005decoding,donoho2006compressed,candestao2006near},
and that earlier underpinned prediction-API model
stealing~\citep{tramer2016stealing}:

\begin{enumerate}[label=\textbf{(P\arabic*)},leftmargin=*]
  \item \textbf{Low-rank codomain.} The logits factor as $z = W_U\,h(c)$ with
  hidden state $h(c)\in\RR^{d}$ and un-embedding $W_U\in\RR^{\Voc\times d}$,
  $d \ll |\Voc|$~\citep{vaswani2017attention}. Hence $\mathrm{range}(f_\theta)$
  lies in a $d$-dimensional subspace, and since $\mathrm{softmax}(z+c\mathbf 1)=\mathrm{softmax}(z)$
  the next-token distribution identifies $z$ only modulo the constant direction
  $\mathrm{span}\{\mathbf 1\}$: the effective output dimension is $d$ (at most $d-1$ for the
  distribution), not $|\Voc|$ --- the softmax bottleneck~\citep{yang2018softmax}.
  \item \textbf{Continuous factorisation.} $f_\theta$ factors through token
  embeddings $c \mapsto E(c)\in\RR^{L\times d}$, so it is the discrete lift of a
  smooth map on a continuous, low-dimensional domain rather than an arbitrary
  function on a combinatorial set.
  \item \textbf{Lipschitz smoothness.} On the compact embedding domain (bounded
  token embeddings plus positional encodings), $f_\theta$ is a composition of
  $C^1$ maps and is therefore Lipschitz, with constant controlled by the product
  of per-layer operator norms. The restriction to a bounded domain is essential:
  standard dot-product self-attention is provably \emph{not} Lipschitz on the
  unbounded domain $\RR^{L\times d}$~\citep[Thm.~3.1]{kim2021lipschitz}, but on a
  ball of radius $R$ its local Lipschitz constant scales as
  $\sqrt{L}\cdot\mathrm{poly}(R,\|W\|)$~\citep[Thm.~3.3]{castin2024smooth} (see
  also \citealp{vuckovic2021regularity}). Thus nearby contexts produce nearby
  logits --- the condition under which point evaluations are informative
  (Section~\ref{sec:rip}).
  \item \textbf{Small description complexity.} For a bounded-weight, fixed
  architecture the metric entropy of the realizable class grows only
  polynomially in $1/\varepsilon$ --- roughly
  $\Hent_\varepsilon(\Fclass)\lesssim W\log(1/\varepsilon)$ with $W$ the parameter
  count --- rather than the $(1/\varepsilon)^{D}$ entropy of generic Lipschitz
  functions on $[0,1]^{D}$, which is exponential in the ambient dimension
  $D$~\citep{kolmogorov1959entropy,anthonybartlett1999}. For the transformer
  class specifically, norm- and rank-based covering-number bounds make the
  dependence on embedding dimension and context length only
  logarithmic~\citep{edelman2022inductive,trauger2024sequence}; networks remain
  Kolmogorov--Donoho rate-distortion optimal approximants of structured function
  classes~\citep{elbraechter2021dnn}.
\end{enumerate}

These four properties are the compressibility priors. The thesis makes precise
how each one reduces the number of evaluations needed.

% =============================================================================
\section{Measurement Model}

Black-box access furnishes an \emph{evaluation operator}
$\Phi:\Fclass\to(\RR^{\Voc})^m$, $\Phi(f)=(f(c_1),\dots,f(c_m))$, instantiated at
several realism levels:

\begin{itemize}[leftmargin=*]
  \item \textbf{Full-logit oracle:} $y_i=f_\theta(c_i)\in\RR^{\Voc}$ (clean linear
  measurements of the codomain).
  \item \textbf{Top-$k$ log-probability oracle:} only the $k$ largest entries are
  observed --- a \emph{compressed/partial} measurement of each output vector.
  \item \textbf{Sampled-token oracle:} only stochastic draws from
  $\mathrm{softmax}(z)$ --- \emph{noisy, quantised} measurements requiring
  repetition to denoise.
  \item \textbf{Adaptive logit-bias oracle:} query-time perturbations
  (\citet{carlini2024stealing}) act as designed measurement vectors that
  reconstruct individual logit coordinates through a top-$k$ API, after which the
  codomain subspace is identified by SVD --- subspace probing under restricted
  output access.
\end{itemize}

A central practical theme is how these restricted oracles degrade the recovery
guarantees of the full-logit case --- i.e.\ how API defenses act as perturbations
of the measurement operator.

% =============================================================================
\section{Core Theoretical Object: Restricted Isometry over $\Fclass$ via Point Evaluation}
\label{sec:rip}

Fix a context distribution $\mathcal{D}$ and the population seminorm
$\|f\|^2=\EE_{c\sim\mathcal{D}}\|f(c)\|^2$. Draw probe contexts
$c_1,\dots,c_m\sim\mathcal{D}$ and form the empirical evaluation seminorm
$\widehat{\|f\|}^2=\tfrac1m\sum_i\|f(c_i)\|^2$. The recovery problem is well-posed
precisely when the evaluation operator is a restricted isometry over the
difference set $\Fclass-\Fclass$.

\begin{condition}[Evaluation-RIP over $\Fclass$]\label{cond:rip}
There exists $\delta\in(0,1)$ such that, for all $f,g\in\Fclass$,
\[
  (1-\delta)\,\|f-g\|^2
  \;\le\;
  \frac1m\sum_{i=1}^{m}\|f(c_i)-g(c_i)\|^2
  \;\le\;
  (1+\delta)\,\|f-g\|^2 .
\]
\end{condition}

Condition~\ref{cond:rip} says: any two models agreeing on the $m$ probe contexts
must be globally close. When it holds with small $\delta$, evaluations at
$\{c_i\}$ form a norm-preserving sketch of the entire function class, and recovery
(find $\hat f\in\Fclass$ consistent with the measurements) is stable. The thesis
develops this object along three established lines.

\paragraph{Sample complexity via metric entropy.} The uniform isometry is an
empirical-process statement; the required $m$ is governed by the metric entropy
$\Hent_\varepsilon(\Fclass)$ of the transformer class. Norm- and rank-based
covering-number bounds for self-attention show $\Hent_\varepsilon(\Fclass)$ scales
polynomially in $1/\varepsilon$ and only logarithmically (or not at all) in the
sequence length~\citep{edelman2022inductive,trauger2024sequence}, so heuristically
$m\gtrsim\Hent_\varepsilon(\Fclass)/(1-\delta)$. The closest prior on recovering a
\emph{system} from input--output traces in a metric-entropy-optimal manner is
\citet{hutter2022metric}.

\paragraph{Well-posedness via smoothness (sampling numbers).} Whether point
samples suffice --- rather than arbitrary linear functionals --- is the subject of
information-based complexity. Recent results show that whenever the approximation
numbers are square-summable (decay rate $s>\tfrac12$), the sampling numbers decay
at the same polynomial rate, so point evaluations are essentially as powerful as
optimal linear measurements~\citep{krieg2021function,krieg2021functionII,novak2010tractability};
the sharp, logarithm-free form $g_{cm}^2\lesssim\tfrac1m\sum_{k\ge m}a_k^2$ is due to
\citet{dolbeault2023sharp}.
The Lipschitz prior (P3) is exactly what places $\Fclass$ in this regime,
justifying that query access loses little against an idealised measurement
operator. Note that these results establish the \emph{rate} of an optimal sampling
algorithm; deriving the two-sided isometry of Condition~\ref{cond:rip} still
requires an empirical-process argument over $\Fclass-\Fclass$.

\paragraph{Lower bounds via packing.} Applying Fano's inequality to a maximal
$\varepsilon$-packing of $\Fclass$ converts metric entropy into a minimax lower
bound: the achievable accuracy $\varepsilon_m$ is fixed by the entropy-balance
$\Hent_{\varepsilon_m}(\Fclass)\asymp m\,\varepsilon_m^2$ for the noisy
(sampled-token) oracle~\citep{yangbarron1999minimax}, so no recovery scheme ---
adaptive or not --- beats this rate. This pins query complexity to the
\emph{intrinsic} complexity of the model class, independent of vocabulary size and
context length except through $\Hent_\varepsilon$.

% =============================================================================
\section{Recoverable Targets}
\label{sec:targets}

Recovering all of $f_\theta$ to high accuracy is infeasible when
$\Hent_\varepsilon(\Fclass)$ is large; the tractable regime is recovery of
low-complexity \emph{targets} derived from $f_\theta$, all under the single
evaluation operator of Section~3.

\begin{target}[Codomain structure]
The subspace $\mathrm{range}(W_U)$ and the hidden dimension $d$ --- the latter being
identifiable, since $W_U$ itself is recoverable only up to an $\Voc$-irrelevant
$d\times d$ symmetry (affine in practice, orthogonal in principle). This is a
functional of $f_\theta$, recoverable as numerical-rank / subspace recovery (SVD of
stacked logit vectors) from $O(d)$ full-logit queries; under a restricted
top-$k$/logit-bias API the cost rises to $O(d\,|\Voc|)$, since each full logit
vector must first be reconstructed coordinate-by-coordinate.
\citet{carlini2024stealing} is the empirical existence proof and the controlled
benchmark.
\end{target}

\begin{target}[Restriction to a structured sub-domain]
$f_\theta|_{\mathcal{S}}$, where $\mathcal{S}$ is a prompt family varying along
$k\ll Ld$ axes (e.g.\ in-context demonstrations of a function
class~\citep{garg2022what}). Recovery is local nonlinear function
approximation with sample complexity controlled by
$\Hent_\varepsilon(\Fclass|_{\mathcal{S}})$, small when $\mathcal{S}$ is
low-dimensional. This subsumes ``identify the algorithm implemented in-context''
as recovery of $f_\theta$ on a structured slice.
\end{target}

\begin{target}[General functionals]
$P(f_\theta)$ --- a decision boundary on a sub-domain, the presence of a
watermark, a behavioural property. Such a functional is recoverable under a
restricted isometry on the relevant quotient of $\Fclass$, typically with far
fewer queries than full recovery.
\end{target}

% =============================================================================
\section{Research Questions}

\begin{description}[leftmargin=*,style=nextline]
  \item[RQ1.] When does the point-evaluation operator satisfy a restricted
  isometry over the transformer-realizable class $\Fclass$, and how do the four
  compressibility priors (P1--P4) enter the isometry constant?
  \item[RQ2.] What is the query complexity $m$ for each target in terms of the
  metric entropy / sampling numbers of the relevant (restricted or quotiented)
  class, and does it match the packing lower bound?
  \item[RQ3.] How do the restricted oracles of Section~3 --- top-$k$ truncation,
  sampling noise, finite precision, rate limits --- degrade the isometry constant,
  and which defenses provably destroy recoverability versus merely inflating $m$?
\end{description}

% =============================================================================
\section{Methodology}

\paragraph{Theory track.} Establish the evaluation-RIP condition for $\Fclass$;
import metric-entropy bounds for transformer-type classes and sampling-number
results to convert them into query-complexity upper bounds; prove matching packing
lower bounds; analyse restricted oracles as perturbations of the isometry.

\paragraph{Empirical track.} On open-weight models where $f_\theta$ is fully
known (Pythia, Qwen, Llama across scales): (i)~empirically estimate the
evaluation-isometry constant for different probe designs (random vs structured vs
adaptively chosen contexts); (ii)~recover targets T1--T3 and compare achieved
query counts against the theory; (iii)~stress-test under simulated top-$k$ and
sampling oracles. Deliverables include a reproducible probing/recovery harness and
empirical query-complexity scaling curves.

% =============================================================================
\section{Evaluation}

\begin{itemize}[leftmargin=*]
  \item \textbf{Recovery accuracy:} subspace/parameter error (T1),
  function-approximation error on $\mathcal{S}$ (T2), functional error (T3).
  \item \textbf{Query efficiency:} measured $m$ vs predicted complexity; scaling
  against $d$, $|\Voc|$, and the dimensionality of $\mathcal{S}$.
  \item \textbf{Isometry diagnostics:} empirical $\delta$ as a function of probe
  design and budget.
  \item \textbf{Robustness:} degradation across the oracle hierarchy; lower-bound
  tightness (achieved vs information-theoretic $m$).
\end{itemize}

% =============================================================================
\section{Expected Contributions}

\begin{enumerate}[leftmargin=*]
  \item A single compressed-sensing formalism for black-box LLMs: the model as one
  nonlinear function, query access as point evaluation, recovery governed by a
  restricted isometry over the function class.
  \item Query-complexity upper bounds (via metric entropy and sampling numbers)
  and matching packing lower bounds for codomain structure, sub-domain
  restrictions, and functionals.
  \item A robustness theory relating realistic API oracles to the isometry
  constant, with provable separations between benign and recovery-destroying
  defenses.
  \item An open, reproducible harness and empirical scaling laws relating query
  budget to recoverable structure.
\end{enumerate}

The framework quantifies \emph{how auditable} a deployed model is from the
outside, and conversely which deployment choices provably limit external recovery
--- a direct contribution to trustworthy-AI and model-auditing goals.

% =============================================================================
\section{Risks and Contingencies}

\begin{itemize}[leftmargin=*]
  \item \emph{The point-evaluation operator may resist a clean RIP for the full
  class.} Mitigation: smoothness places $\Fclass$ in the sampling-numbers regime;
  where a global RIP is unavailable, restrict attention to targets T1--T3, whose
  quotient classes have small metric entropy and admit isometries directly.
  \item \emph{Restricted oracles may block logit access entirely.} Mitigation: the
  open-weight track is independent of any external API; oracle restrictions are
  simulated locally with controlled corruption layers.
  \item \emph{Theory--experiment gap.} Mitigation: T1 (validated by
  \citet{carlini2024stealing}) guarantees a defensible empirical result even if
  the bounds for T2/T3 come out only partial.
\end{itemize}

% =============================================================================
\section{Indicative Timeline (6 months)}

\begin{center}
\begin{tabular}{cl}
\toprule
\textbf{Month} & \textbf{Milestone}\\
\midrule
1 & Formalise the function/measurement model; consolidate metric-entropy and sampling-number tools\\
2 & Evaluation-RIP theory; codomain recovery (T1) + benchmark harness\\
3 & T1 empirical validation; oracle-robustness study\\
4 & Sub-domain restriction (T2) and functionals (T3): complexity upper/lower bounds\\
5 & T2/T3 experiments; query-complexity scaling laws\\
6 & Synthesis, writing, defense preparation\\
\bottomrule
\end{tabular}
\end{center}

% =============================================================================
\section{Prerequisites for the Student}

Linear algebra and high-dimensional probability; prior exposure to compressed
sensing or statistical learning theory (sample complexity, covering/packing
arguments); comfort with PyTorch and running open-weight LLMs. Familiarity with
metric-entropy / information-based-complexity arguments is an asset but can be
acquired during Month~1.

% =============================================================================
\bibliographystyle{unsrtnat}
\bibliography{references}

\end{document}
